% !TeX root = ../../Book3.tex
% 纲要标题：### 21.2 Hilbert 合冲定理
% 纲要位置：计划纲要.md，第 2629 行。

\section{Hilbert 合冲定理}
\label{b3:sec:2102}


% 纲要内容（仅作写作计划，尚非教材正文）：
%
% * 添加变量的两种作用之差；通过张量投射消解、提升链映射及映射锥给出整体维数上界
% * 区分全部模的投射维数上界与有限生成分次模的有限自由消解；既有分次下的投射模自由性
% * 自由消解长度与变量数；最高合冲模与截短消解的指标
%
% ---
%

Noether 性保证每一层关系有限生成，尚不能阻止关系层数无限增加。多项式环的特殊之处是变量数给出统一上界。先证明对任意模成立的投射维数上界，再用分次极小性把它转为有限自由消解；这样无须提前假设一般有限投射模都自由。

\subsection{多项式环的有限整体维数}
\label{b3:subsec:2102:01}

以下整体维数取所有模的投射维数的上确界，包括无限生成模。增加一个变量时，有一条简单的正合列比较新旧环上的模。

\ProseParagraphBreak
设 \(B=A[t]\)、\(M\) 为 \(B\) 模。\(B\otimes_AM\) 中暂时有两种 \(t\) 作用，记为
\[
t_{\mathrm{poly}}(b\otimes m)=bt\otimes m,
\qquad t_M(b\otimes m)=b\otimes tm.
\]
前者来自第一因子上新添的变量，后者来自 \(M\) 原有的作用。要恢复 \(M\)，必须让这两个作用一致，所以取 \(\delta=t_{\mathrm{poly}}-t_M\) 作为关系映射。例如 \(A=k\)、\(M=k\)，且 \(t\) 在 \(M\) 上按标量 \(\lambda\in k\) 作用，则 \(B\otimes_kM=k[t]\)，相应序列就是
\[
0\longrightarrow k[t]\xrightarrow{\,t-\lambda\,}k[t]
\xrightarrow{f\mapsto f(\lambda)}k\longrightarrow0.
\]
多出的变量带来一层使两种作用一致的关系；下列引理证明同一构造对任意 \(A\) 和 \(M\) 仍成立。

\begin{lemma}\label{b3:lem:polynomial-module-sequence}
设 \(A\) 为交换环，\(B=A[t]\)，\(M\) 为任意 \(B\) 模。把 \(M\) 限制为 \(A\) 模后，有自然的 \(B\) 模短正合列
\[
0\longrightarrow B\otimes_AM\xrightarrow{\delta}B\otimes_AM
\xrightarrow{\varepsilon}M\longrightarrow0,
\]
其中 \(\delta(b\otimes m)=bt\otimes m-b\otimes tm\)，\(\varepsilon(b\otimes m)=bm\)。
\end{lemma}
\begin{proof}
两式尊重 \(A\) 平衡关系并且 \(B\) 线性，且 \(\varepsilon\delta=0\)、\(\varepsilon(1\otimes m)=m\)。作为 \(A\) 模，\(B\otimes_AM\) 是各 \(t^i\otimes M\) 的直和，每个元素只有有限个非零分量。若 \(u=\sum_{i=0}^rt^i\otimes m_i\ne0\)、\(m_r\ne0\)，则 \(\delta(u)\) 的最高 \(t\) 次分量为 \(t^{r+1}\otimes m_r\)，故 \(\delta\) 单射。这里比较的只是 \(t\) 的指数，没有使用单项式序。模掉其像，可以反复用 \(t^i\otimes m\equiv t^{i-1}\otimes tm\)，最终得
\(u\equiv1\otimes\sum_i t^im_i=1\otimes\varepsilon(u)\)。所以 \(\ker\varepsilon=\operatorname{im}\delta\)。
\end{proof}

\begin{proposition}\label{b3:prop:polynomial-global-bound}
设 \(A\) 为交换环，\(d\ge0\) 为整数，且 \(A\) 的整体维数至多为 \(d\)。则 \(A[t]\) 的整体维数至多为 \(d+1\)。因而对任意域 \(k\) 与非负整数 \(n\)，多项式环 \(k[x_1,\ldots,x_n]\) 的整体维数至多为 \(n\)。
\end{proposition}
\begin{proof}
置 \(B=A[t]\)，任取一个 \(B\) 模 \(M\)。上一引理把 \(M\) 表示为两份 \(T=B\otimes_AM\) 之间单射 \(\delta\) 的余核。先从 \(M\) 的 \(A\) 投射消解扩标量，得到这两份 \(T\) 的 \(B\) 投射消解，再将 \(\delta\) 提升为链映射。映射锥把两条消解合成一条复形，其中一份向高一层平移，正是长度上界增加一的原因。

取 \(M\) 作为 \(A\) 模的长度至多 \(d\) 的投射消解。\(B\) 在 \(A\) 上自由，因而平坦，故张量 \(B\) 后消解仍正合。它的各项也是投射 \(B\) 模：若投射 \(A\) 模 \(P\) 满足 \(P\oplus P''\simeq A^{(I)}\)，则
\[
(B\otimes_AP)\oplus(B\otimes_AP'')\simeq B^{(I)},
\]
所以 \(B\otimes_AP\) 是自由 \(B\) 模的直和项。

记 \(T=B\otimes_AM\)，在上一引理的两份 \(T\) 上取这样得到的投射消解 \(P'_\bullet,Q'_\bullet\)，增广分别为 \(\varepsilon_P:P'_0\rightarrow T\)、\(\varepsilon_Q:Q'_0\rightarrow T\)。由 \(P'_0\) 投射和 \(\varepsilon_Q\) 满射，可选 \(a_0:P'_0\rightarrow Q'_0\)，使
\(\varepsilon_Qa_0=\delta\varepsilon_P\)。在第 \(1\) 层，\(a_0\mathrm{d}_P(P'_1)\subseteq\ker\varepsilon_Q\)；由 \(Q'_1\rightarrow\ker\varepsilon_Q\) 满射和 \(P'_1\) 投射，可将 \(a_0\mathrm{d}_P\) 提升为 \(a_1\)。以后若已构造到第 \(j-1\) 层，链映射等式使 \(a_{j-1}\mathrm{d}_P(P'_j)\) 落入 \(Q'_{j-1}\) 的微分核，正合性与 \(P'_j\) 的投射性便给出下一次提升。这样得到覆盖 \(\delta\) 的链映射 \(a:P'_\bullet\rightarrow Q'_\bullet\)。

要把模层的余核转成消解，可以让一项同时包含目标消解的第 \(i\) 层与源消解的第 \(i-1\) 层，利用 \(a\) 将两者连接。删去两条消解的增广，约定负次数项为零，构造链复形
\[
C_i=Q'_i\oplus P'_{i-1},\qquad
D_i(q,p)=(\mathrm{d}_Qq+a_{i-1}p,-\mathrm{d}_Pp)\quad(i\ge1)
\]
以及增广 \(C_0=Q'_0\xrightarrow{\varepsilon_Q}T\xrightarrow{\varepsilon}M\)。链映射等式给出 \(D^2=0\)，且增广与 \(D_1\) 的复合为零。这个复形称为 \(a\) 的映射锥，其各项投射。两条原消解满足 \(P'_i=Q'_i=0\) 对所有 \(i>d\) 成立，因此
\[
C_i=Q'_i\oplus P'_{i-1}=0\qquad(i>d+1).
\]
最高可能非零的一层是 \(C_{d+1}=P'_d\)。接下来检验同调：映射锥将 \(\delta\) 的核放在一次同调、余核放在零次同调；本例中 \(\delta\) 单射，因而只留下所需的余核 \(M\)。

短正合复形
\[
0\longrightarrow Q'_\bullet\longrightarrow C_\bullet
\longrightarrow P'_\bullet[1]\longrightarrow0,
\qquad P'[1]_i=P'_{i-1},\quad
\mathrm{d}_{P'[1]}=-\mathrm{d}_P
\]
按\cref{b3:lem:tor-basic-properties}证明中的圈提升构造给出同调长正合列。因为 \(P'\) 与 \(Q'\) 的正次同调为零，\(H_i(C)=0\) 对 \(i\ge2\) 成立；其余部分为
\[
0\longrightarrow H_1(C)\longrightarrow H_0(P')
\xrightarrow{H_0(a)}H_0(Q')\longrightarrow H_0(C)
\longrightarrow0.
\]
在 \(H_0(P')\simeq T\simeq H_0(Q')\) 的识别下，\(H_0(a)=\delta\)：圈 \(p\in P'_0\) 在 \(C_1\) 中的提升为 \((0,p)\)，其微分是 \((a_0p,0)\)。上一引理已证明 \(\delta\) 单射且余核为 \(M\)，所以 \(H_1(C)=0\)、\(H_0(C)\simeq M\)，增广后的 \(C\) 是所需的投射消解。于是任意 \(M\) 的投射维数至多 \(d+1\)。域上所有向量空间自由，整体维数为零，逐个添加变量完成归纳。
\end{proof}

\subsection{Hilbert 合冲定理及其证明}
\label{b3:subsec:2102:02}

\begin{theorem}[Hilbert 合冲定理]\label{b3:thm:hilbert-syzygy}
设 \(k\) 为域，\(S=k[x_1,\ldots,x_n]\) 取标准分次。每个有限生成分次 \(S\) 模都有长度至多为 \(n\)、各项均为有限秩的分次自由消解；它的分次极小自由消解在第 \(n\) 层以后全部为零。

此外，\(S\) 的整体维数恰为 \(n\)：任意 \(S\) 模的投射维数都不超过 \(n\)，这里不要求有限生成或分次，而 \(k=S/(x_1,\ldots,x_n)\) 的投射维数恰为 \(n\)。
\end{theorem}
\begin{proof}
由\cref{b3:prop:graded-minimal-existence}先构造可能无限的极小消解 \(F_\bullet\)。上一命题说明 \(M\) 有长度至多 \(n\) 的投射消解，故 \(i>n\) 时 \(\operatorname{Tor}_i^S(M,k)=0\)。由式\eqref{b3:eq:graded-betti-tor}得 \(F_i\otimes_Sk=0\)，所以 \(F_i=0\)，得到所需有限自由消解。

另一方面，\(x_1,\ldots,x_n\) 是 \(S\) 正则序列：逐次商为剩余变量的多项式环，每一步所乘变量都是非零因子，末端商为 \(k\ne0\)。\Cref{b3:thm:koszul-regular-resolution}因此给出 \(k\) 的自由消解。张量 \(k\) 后微分为零，第 \(n\) 项为 \(\bigwedge^n k^n\simeq k\)，故 \(\operatorname{Tor}_n^S(k,k)\ne0\)。这迫使 \(\operatorname{pd}_Sk\ge n\)，结合上界得 \(\operatorname{gldim}S=n\)。\(n=0\) 时环为域，同样成立。
\end{proof}

这就是\term[Hilbert-hechong-dingli]{Hilbert 合冲定理}[Hilbert's syzygy theorem]的分次形式。证明中有两项不同的有限性：有限生成性使每一项自由模具有有限秩，变量数则使非零项的数目有统一上界。整体维数的结论取遍任意 \(S\) 模，给出的是投射消解的长度界。对有限生成无分次模，若要把有限生成投射项进一步全局自由化，还需要关于多项式环上有限生成投射模的额外定理；本章关于自由性的论证使用既有分次与极小消解。

\subsection{自由消解长度与变量数}
\label{b3:subsec:2102:03}

上界可以达到：\(k=S/(x_1,\ldots,x_n)\) 的 Koszul 消解长度恰为 \(n\)，其 Betti 数为 \(\beta_{i,i}=\binom ni\)。若只商掉前 \(r\) 个变量，则相同证明给出长度恰为 \(r\) 的极小消解。相反，一个自由模的消解长度为零，与所选生成元次数无关。

\ProseParagraphBreak
长度也可以从合冲模辨认。设 \(k\) 为域，\(S=k[x_1,\ldots,x_n]\) 取标准分次，\(0\ne M\) 为有限生成分次 \(S\) 模，其极小自由消解满足 \(F_r\ne0\)、\(F_{r+1}=0\)。于是 \(\operatorname{Tor}_r^S(M,k)\ne0\)，所以不存在更短的投射消解。若 \(r=0\)，则增广给出 \(F_0\simeq M=\Omega^0M\)。若 \(r\ge1\)，记 \(\mathrm{d}_0=\varepsilon\)；正合性及 \(\mathrm{d}_r\) 单射给出
\[
\Omega^rM=\ker\mathrm{d}_{r-1}
=\operatorname{im}\mathrm{d}_r\simeq F_r.
\]
这里 \(\Omega^rM\) 是 \(F_{r-1}\) 的子模，\(F_r\) 经单射 \(\mathrm{d}_r\) 同构映到它；三个下标分别记录合冲层数、核所在的前一层和产生这个像的自由项。
如果某个 \(0\le j<r\) 的 \(\Omega^jM\) 已是投射模，则 \(j=0\) 时 \(M\) 本身投射，\(j\ge1\) 时可截得投射消解
\[
0\longrightarrow\Omega^jM\longrightarrow F_{j-1}
\longrightarrow\cdots\longrightarrow F_0
\longrightarrow M\longrightarrow0.
\]
两种情形都将消解长度缩至 \(j\)，与 \(\operatorname{Tor}_r^S(M,k)\ne0\) 矛盾。因而“已经找到某个有限消解”与“已经找到最短长度”是两种结论；后者需要极小性或非消失同调作为证据。

\ProseParagraphBreak
例如 \(S=k[x,y]\)、\(M=S/(x^2,xy)\) 有消解
\[
0\longrightarrow S(-3)\xrightarrow{\binom{y}{-x}}S(-2)^2
\xrightarrow{(x^2\ xy)}S\longrightarrow M\longrightarrow0.
\]
正合性由 \(ax^2+bxy=0\) 在整环中约去 \(x\)，再使用 \(ax+by=0\) 的计算得到；全部非零微分项在 \((x,y)\) 中，故消解极小，投射维数为二。商环的维数却是一，因为 \(\sqrt{(x^2,xy)}=(x)\)。为使用 Auslander--Buchsbaum 公式，置 \(\mathfrak m=(x,y)\)，局部化到正则局部环 \(S_{\mathfrak m}\)。上述消解局部化后仍极小，各非零自由项仍非零，故 \(\operatorname{pd}_{S_{\mathfrak m}}M_{\mathfrak m}=2\)。\Cref{b3:thm:auslander-buchsbaum}给出
\[
\operatorname{depth}_{S_{\mathfrak m}}M_{\mathfrak m}
=\operatorname{depth}S_{\mathfrak m}-2=0.
\]
在这个正则局部环境中，长度二记录的是环境深度与模深度之差；只在 CM 情形，才由\cref{b3:prop:cm-pd-codimension}把它解释为支集余维。

\begin{proposition}\label{b3:prop:graded-projective-free}
设 \(k\) 为域，\(S=k[x_1,\ldots,x_n]\) 取标准分次，\(P\) 为有限生成分次 \(S\) 模。若 \(P\) 的底层 \(S\) 模投射，则 \(P\) 同构于有限个 \(S(-a)\) 的直和，即为有限秩分次自由模。
\end{proposition}
\begin{proof}
置 \(\mathfrak m=(x_1,\ldots,x_n)\)，取 \(P/\mathfrak mP\) 的齐次基并提升，得到分次极小满射 \(F\to P\)，其中 \(F\) 为有限秩分次自由模，核 \(K\subseteq\mathfrak mF\) 齐次且有限生成。由于底层 \(P\) 投射，短正合列
\[
0\longrightarrow K\longrightarrow F\longrightarrow P\longrightarrow0
\]
在底层模范畴中分裂，所以张量 \(k=S/\mathfrak m\) 后仍正合。原映射诱导的 \(F/\mathfrak mF\to P/\mathfrak mP\) 又是同构，因而 \(K/\mathfrak mK=0\)。\Cref{b3:lem:graded-nakayama}使 \(K=0\)，故原分次满射为同构，给出所需的齐次自由基。这里使用了 \(P\) 的既有分次，底层分裂本身无需预先保持次数。
\end{proof}
